% Online appendix body. Generated by write_appendix_body.py from appendix_template.tex and output/appendix.json; edit the template.
\section*{Online appendix}
\setcounter{table}{0}
\renewcommand{\thetable}{A\arabic{table}}
\setcounter{section}{0}
\renewcommand{\thesection}{A.\arabic{section}}

\section{Data, audit and sample construction}
The Greater Austin panel is read from the FCACTX survey publications for 2013--2024 (there was no 2025 survey; the 2026 wave is not used). Prices and effective dates are taken from the published tables with layout-preserving rules, and unresolved values stay missing. Establishments are linked across waves by telephone number. A two-year link compares survey $t-2$ with survey $t$ for the same establishment, whether or not $t-1$ is observed, using the same rules as the annual links.

The 45 intervals initially flagged as same-date repricing were checked against the current- and preceding-year survey pages. Four price-row alignment errors and one date-parsing error removed five false flags; one real change became ineligible under the two-comparable-measure rule; and correcting a preceding date moved one adjacent interval into the same-date category, giving $45-4-1-1+1=40$. Ten randomly selected final flags (seed 20260907) and two targeted cases were then rechecked. All twelve kept their flag; one required a price-level correction (Gabriels Funeral Chapel, 2014). The audit verifies consistency with the survey pages, not with every underlying provider document.

Of 435 date-defined consecutive intervals, 87 lack two comparable headline prices. Setting every missing new-date outcome to no repricing and every missing same-date outcome to repricing gives a lower bound on the gap of $197/300-53/135=0.264$. The bound does not address records without a usable date.

Kansas City comparisons use FCA-GKC's 2018, 2020 and 2023 publications, which report month-level GPL dates. Branches are linked across publications, and records that the 2020 publication marks as retained from 2018 are excluded at either endpoint. Because dates are month--year fields, an unchanged month is not necessarily an unchanged day, and inclusion of the crematory fee in the 2018 direct-cremation measure is not explicitly established.

\section{The date signal by price measure and observation gap}
Table~\ref{tab:A1} scores the date against each price separately. For the basic-services fee, a single GPL line, the date is almost perfectly sensitive (<<FEE_SENS>>): the fee changed without a new date in <<FEE_SAME>> links. Its specificity is low (<<FEE_SPEC>>) because most new lists leave the fee unchanged. The two bundle prices change under an unchanged date more often (<<DC_SAME>> for direct cremation and <<IB_SAME>> for burial). Over two-year links, few dates stay unchanged, so sensitivity rises and specificity falls. The date has also become more informative over time: $J$ is <<J_EARLY>> for intervals ending in 2014--2019 and <<J_LATE>> for 2020--2024.

\input{tables/appA1_measures}

\section{Same-date changes and survey noise}
Table~\ref{tab:A2} compares same-date and new-date price changes on two symptoms of recording noise: decreases, and reversal to the earlier price in the following year. Same-date changes show more of both, and almost all of the excess is in the burial measure, whose least-expensive casket is priced on a separately dated casket price list. Removing every reversed change leaves the result in the main text. If changes whose next-year price is unobserved are also dropped, the rates are <<PD_RATES>>, with $J=<<PD_J>>$. Same-date repricing occurs at <<SD_EST>> establishments, and none accounts for more than <<SD_MAX>>\% of the 40 intervals.

\input{tables/appA2_noise}

\section{Robustness of the repricing gap}
Table~\ref{tab:A3} collects alternative samples, thresholds and exclusions. Except for the worst-case missing-data bound (<<WORST>>), the gap between new-date and same-date repricing stays between <<GAP_MIN>> and <<GAP_MAX>>, and no single brand or establishment drives it.

\input{tables/appA3_robust}

\section{Prior list age: additional specifications and simulation}
Table~\ref{tab:A4} splits the price regression by date status and repeats both regressions on the sample at which direct cremation and the fee are comparable, where the date and price gradients are close. Table~\ref{tab:A5} reports the simulation described in the main text. With establishment effects, list age predicts a date change even though the simulated hazard is constant, because age resets when the date changes and the within transformation uses each provider's own average date-change rate. Without establishment effects, providers that rarely re-date both hold older lists and re-date less, which produces a negative coefficient.

\input{tables/appA4_age}
\input{tables/appA5_sim}
